% Generated draft: requires mathematical and source review.
\documentclass[11pt]{article}
\usepackage[margin=25mm]{geometry}
\usepackage{fontspec}
\setmainfont{Latin Modern Roman}
\usepackage{amsmath,amssymb,mathtools,unicode-math}
\setmathfont{Latin Modern Math}
\usepackage{enumitem,needspace,xurl,hyperref}
\setlength{\parindent}{0pt}
\setlength{\parskip}{5pt}
\setlength{\emergencystretch}{3em}
\newcommand{\N}{\mathbb{N}}
\newcommand{\Z}{\mathbb{Z}}
\newcommand{\Q}{\mathbb{Q}}
\newcommand{\R}{\mathbb{R}}
\newcommand{\C}{\mathbb{C}}
\newcommand{\F}{\mathbb{F}}
\newcommand{\A}{\mathbb{A}}
\newcommand{\PP}{\mathbb{P}}
\newcommand{\E}{\mathbb{E}}
\newcommand{\eps}{\varepsilon}
\newcommand{\dd}{\,\mathrm{d}}
\begin{document}
\section*{UnsolvedMath ID: 800623}
\textbf{Source problem code:} catalogue entry
\subsection{Definitions and mathematical statement}
Equip every subset of \(\R\) with its usual subspace topology. A subset \(A\) of a space \(T\) is nowhere dense in \(T\) if \(\operatorname{int}_T(\overline{A}^{\,T})=\varnothing\). A set \(V\subseteq\R\) is a \(\nu\)-set if every subset of \(V\) that is nowhere dense in \(V\) is countable; the source describes this as being Luzin relative to itself.

Use the standard meaning of concentration: a space \(Y\) is concentrated on \(D\subseteq Y\) if, for every open set \(U\subseteq Y\) containing \(D\), the set \(Y\setminus U\) is countable. Here countable includes finite and empty sets. The supplied record does not explicitly define concentration, so this interpretation requires source confirmation.

The question is whether the following implication holds for every \(X\subseteq\R\):
\[
\begin{gathered}
\forall Y\subseteq X\ \exists D\subseteq Y\quad
\left[ D\text{ is countable and }
\forall U\subseteq Y\text{ open in }Y,
\quad D\subseteq U\Longrightarrow Y\setminus U\text{ is countable}\right]
\\[2pt]
\Longrightarrow
\exists V\subseteq\R\ \exists f:V\to X\quad
\left[V\text{ is a }\nu\text{-set and }f\text{ is continuous and surjective}\right].
\end{gathered}
\]
The supplied source states that every continuous image of a \(\nu\)-set satisfies the hypothesis; the question asks for the converse.
\subsection{Short English statement}
Suppose every subset \(Y\) of a set of reals \(X\) has a countable subset \(D\) such that every open neighborhood of \(D\) in \(Y\) leaves only countably many points of \(Y\) outside it. Must \(X\) be a continuous image of a set of reals whose every relatively nowhere dense subset is countable?
\subsection{Sources}
\begin{itemize}
\item Local UnsolvedMath registry, permanent ID \(800623\), source problem code \texttt{catalogue entry}; supplied upstream revision \texttt{d371ee917ac9647fbb06cd6fc529deec6f3dda82}. The registry labels the problem open; this label is supplied metadata, not an independently verified status assessment.
\item Tsaban, \emph{Selection Principles and Special Sets of Reals: Open Problems} (2006), Problem 52, \url{https://arxiv.org/abs/math/0609601}. The record reports access on 2026-10-16 and extraction from arXiv LaTeX; these provenance details are supplied and have not been independently verified here.
\item The supplied attribution/context is \texttt{[BrownCox83]}. No full bibliographic entry for this reference is supplied. The registry lists Tsaban as proposer and 2006 as the proposed year; the relationship between this metadata and the earlier attribution is unspecified.
\end{itemize}
\end{document}
